%==============================================================================
%  CARA --- Project Book
%  Final Degree Project, B.Sc. in Computer Science, Deep Learning Major
%
%  Compile with: pdfLaTeX  (Overleaf: Menu → Compiler → pdfLaTeX)
%  Bibliography: BibTeX (biblatex + biber also works if preferred)
%
%  This file is self-contained: all diagrams are drawn with TikZ so the
%  document compiles on Overleaf without uploading any external images.
%  Replace the TikZ diagrams with \includegraphics{...} once you have
%  screenshots / figures ready.
%==============================================================================
\documentclass[11pt,a4paper,oneside]{report}

%------------------------------------------------------------------------------
% Encoding & fonts
%------------------------------------------------------------------------------
\usepackage[T1]{fontenc}
\usepackage[utf8]{inputenc}
\usepackage{lmodern}
\usepackage{microtype}
\usepackage{textcomp}

%------------------------------------------------------------------------------
% Page layout
%------------------------------------------------------------------------------
\usepackage[a4paper,left=2.8cm,right=2.8cm,top=2.8cm,bottom=3cm]{geometry}
\usepackage{setspace}
\onehalfspacing
\usepackage{parskip}

%------------------------------------------------------------------------------
% Math, symbols
%------------------------------------------------------------------------------
\usepackage{amsmath,amssymb,amsthm,mathtools}
\usepackage{siunitx}
\sisetup{detect-weight=true,detect-family=true}

%------------------------------------------------------------------------------
% Graphics & diagrams
%------------------------------------------------------------------------------
\usepackage{graphicx}
\usepackage{xcolor}
\usepackage{tikz}
\usetikzlibrary{arrows.meta,positioning,shapes.geometric,fit,calc,shadows.blur,backgrounds,decorations.pathreplacing}
\usepackage{pgfplots}
\pgfplotsset{compat=1.18}

%------------------------------------------------------------------------------
% Tables
%------------------------------------------------------------------------------
\usepackage{booktabs}
\usepackage{array}
\usepackage{tabularx}
\usepackage{multirow}
\usepackage{longtable}
\usepackage{caption}
\captionsetup{font=small,labelfont=bf,skip=6pt}
\usepackage{subcaption}

%------------------------------------------------------------------------------
% Code listings
%------------------------------------------------------------------------------
\usepackage{listings}
\definecolor{codebg}{RGB}{248,248,248}
\definecolor{codekey}{RGB}{0,90,160}
\definecolor{codestr}{RGB}{163,21,21}
\definecolor{codecmt}{RGB}{80,120,80}
\lstdefinestyle{cara}{
    backgroundcolor=\color{codebg},
    basicstyle=\ttfamily\small,
    keywordstyle=\color{codekey}\bfseries,
    stringstyle=\color{codestr},
    commentstyle=\color{codecmt}\itshape,
    numbers=left,
    numberstyle=\tiny\color{gray},
    numbersep=8pt,
    frame=single,
    rulecolor=\color{gray!30},
    breaklines=true,
    breakatwhitespace=true,
    showstringspaces=false,
    tabsize=2,
    captionpos=b,
    xleftmargin=12pt,
    framexleftmargin=8pt,
}
\lstset{style=cara}

%------------------------------------------------------------------------------
% Chapter / section styling
%------------------------------------------------------------------------------
\usepackage{titlesec}
\definecolor{caraprimary}{RGB}{30,60,120}
\definecolor{caraaccent}{RGB}{200,80,60}
\titleformat{\chapter}[display]
    {\normalfont\huge\bfseries\color{caraprimary}}
    {\filleft\Large\textcolor{gray!70}{Chapter \thechapter}}
    {0.6ex}
    {\Huge}
    [\vspace{-0.4ex}\rule{\linewidth}{0.6pt}]
\titlespacing*{\chapter}{0pt}{-30pt}{22pt}

\titleformat{\section}
    {\normalfont\Large\bfseries\color{caraprimary}}{\thesection}{1em}{}
\titleformat{\subsection}
    {\normalfont\large\bfseries\color{caraprimary!85}}{\thesubsection}{1em}{}

%------------------------------------------------------------------------------
% Headers / footers
%------------------------------------------------------------------------------
\usepackage{fancyhdr}
\pagestyle{fancy}
\fancyhf{}
\renewcommand{\headrulewidth}{0.4pt}
\fancyhead[L]{\small\textsc{CARA --- Project Book}}
\fancyhead[R]{\small\leftmark}
\fancyfoot[C]{\thepage}
\renewcommand{\chaptermark}[1]{\markboth{\thechapter.\ #1}{}}

%------------------------------------------------------------------------------
% Hyperlinks (load late)
%------------------------------------------------------------------------------
\usepackage[hidelinks,bookmarksnumbered,colorlinks=true,linkcolor=caraprimary,citecolor=caraaccent,urlcolor=caraprimary]{hyperref}
\usepackage{cleveref}

%------------------------------------------------------------------------------
% Convenience macros
%------------------------------------------------------------------------------
\newcommand{\cara}{\textsc{Cara}\xspace}
\usepackage{xspace}
\newcommand{\effnet}{EfficientNet-B0\xspace}
\newcommand{\pytorch}{PyTorch\xspace}

\newtheorem{observation}{Observation}[chapter]

%==============================================================================
%  COVER PAGE
%==============================================================================
\begin{document}
\pagenumbering{roman}

\begin{titlepage}
    \centering
    \vspace*{0.5cm}
    {\scshape\large \textcolor{caraprimary}{[UNIVERSITY / COLLEGE NAME]}\par}
    {\scshape\normalsize [Faculty of Computer Science] \quad --- \quad [Deep Learning Major]\par}
    \vspace{2.5cm}

    \rule{\linewidth}{1.2pt}\vspace{0.4cm}
    {\Huge\bfseries \textcolor{caraprimary}{CARA}\par}
    \vspace{0.3cm}
    {\LARGE\itshape Computer-Aided Review \& Analysis of Facial Skin\par}
    \vspace{0.4cm}
    {\large A Multi-Model Deep Learning System for Automated Skin Assessment\\ and Evidence-Based Skincare Recommendation\par}
    \vspace{0.2cm}
    \rule{\linewidth}{1.2pt}

    \vspace{2.0cm}
    {\large \textbf{Final Degree Project}\par}
    {\normalsize B.Sc. in Computer Science --- Deep Learning Track\par}

    \vfill
    \begin{tabular}{rl}
        \textbf{Submitted by:} & [Student Name 1] \\
                               & [Student Name 2] \\
                               & [Student Name 3] \\
        \textbf{Supervisor:}   & [Prof.\ / Dr.\ Supervisor Name] \\
        \textbf{Department:}   & [Department Name] \\
        \textbf{Submission:}   & [Month, Year] \\
    \end{tabular}

    \vspace{1.2cm}
    {\footnotesize \textit{CARA is not a medical diagnostic tool.
    It provides cosmetic skincare insights and educational ingredient
    recommendations only.}\par}
\end{titlepage}

%==============================================================================
%  ABSTRACT
%==============================================================================
\chapter*{Abstract}
\addcontentsline{toc}{chapter}{Abstract}

Skincare decisions are today driven mostly by intuition, marketing, and social
media influencers, while professional dermatological advice is expensive and
often inaccessible. \cara is a web application that closes part of this gap by
delivering an honest, evidence-based assessment of a user's facial skin from a
single photograph, together with concrete ingredient-level recommendations.

The system decomposes the assessment into three complementary
classification tasks --- acne severity, skin type, and secondary skin
issues --- each handled by an \effnet convolutional neural network trained
independently. Predictions are combined through a deterministic, rule-based
Business Logic Programming (BLP) engine that turns raw model outputs into
actionable, explainable advice.

We report on the full engineering lifecycle: data collection and cleaning
(built primarily around the ACNE04 dataset, augmented with in-house
photographs), preprocessing (face detection, CLAHE brightness normalisation,
ImageNet normalisation), two-phase transfer-learning training with weighted
loss to address strong class imbalance, and the failed experiments that
shaped our final design (single monolithic classifier, learned late fusion,
naive resizing without face crop).

The final models achieve \SI{69.6}{\percent}, \SI{77.8}{\percent} and
\SI{98.3}{\percent} test accuracy on acne severity, skin type and skin
issues respectively. Beyond raw accuracy, the report focuses on
\emph{truthfulness} --- the guarantee that when the system is uncertain, it
says so --- which we consider a stronger contribution than any single
percentage point of accuracy.

\bigskip
\noindent\textbf{Keywords:} deep learning, computer vision, convolutional
neural networks, EfficientNet, transfer learning, image classification,
dermatology, skincare, rule-based systems, explainable AI, FastAPI, React.

\clearpage

%==============================================================================
%  ACKNOWLEDGEMENTS
%==============================================================================
\chapter*{Acknowledgements}
\addcontentsline{toc}{chapter}{Acknowledgements}

We would like to thank our supervisor, [Supervisor Name], for the guidance,
the patience, and the direct feedback that pushed the project from a demo
into a system we are willing to stand behind. We thank the members of the
[Department Name] for the technical infrastructure and the open-door policy
whenever we hit a wall. Finally, we thank our families and friends for
tolerating months of conversations that revolved around loss curves and
skin types.

\clearpage

%==============================================================================
%  TABLE OF CONTENTS
%==============================================================================
\tableofcontents
\clearpage
\listoffigures
\addcontentsline{toc}{chapter}{List of Figures}
\listoftables
\addcontentsline{toc}{chapter}{List of Tables}
\clearpage

\pagenumbering{arabic}
\setcounter{page}{1}

%==============================================================================
\chapter{Introduction and Motivation}
%==============================================================================

\section{The problem we set out to solve}

Skin is the largest organ in the human body, and yet the decisions people
make about their skin are among the most poorly informed decisions in
everyday life. Consumers spend billions of dollars each year on skincare
products chosen almost entirely on the basis of marketing, celebrity
endorsements, and viral trends. Dermatologists are expensive and
overbooked; general practitioners rarely have time for cosmetic concerns;
and the loudest voices online are often those with the least
qualifications.

The result is a market in which a person with mild acne might spend a year
using a moisturiser designed for dry, ageing skin, while another person
with sensitive combination skin scrubs their face with an abrasive
salicylic-acid cleanser marketed for severe cystic acne. Neither person is
lazy or foolish. They simply lack access to a truthful, personal, and
actionable assessment of their own skin.

\cara --- \emph{Computer-Aided Review \& Analysis} --- is our attempt to
provide a small slice of that missing information. The goal is deliberately
narrow: given a single facial photograph uploaded through a web browser,
return within a few seconds an honest report describing the user's acne
severity, skin type, and dominant secondary issues, together with a
concrete list of ingredients to look for and ingredients to avoid.

\section{Why this problem is a good fit for deep learning}

Dermatological image classification has three properties that make it
particularly interesting from a deep-learning point of view:

\begin{enumerate}
    \item \textbf{Highly visual signal.} Almost everything that a
        dermatologist evaluates in a routine cosmetic consultation is
        visible in a well-lit photograph: comedones, papules, pustules,
        pore density, oily sheen, dry flakes, hyperpigmentation, fine
        lines.
    \item \textbf{Weak, noisy labels.} Public datasets are labelled by
        different clinicians, with different rating scales, in different
        lighting conditions. This forces us to think seriously about
        preprocessing, augmentation, and class imbalance rather than
        assuming clean data.
    \item \textbf{Multiple correlated but distinct sub-tasks.} Acne
        severity, skin type, and skin issues are related --- an oily skin is
        more prone to acne, wrinkles are correlated with skin dryness ---
        but they are not the same task. This creates a natural design
        question: one big model or several small ones?
\end{enumerate}

\section{Non-negotiable design principles}

Very early in the project we agreed on three rules that would govern every
technical decision. We refer back to them throughout the report.

\begin{description}
    \item[Truthfulness is the product.] A wrong prediction is worse than
        no prediction. If the model is not confident, the system must say
        so out loud rather than paper over the uncertainty.
    \item[No manipulated results.] We never tune a threshold to make an
        embarrassing prediction disappear. Numbers reported in this book
        are the numbers we measured.
    \item[Explainability comes from design, not post-hoc excuses.] The
        rule engine that converts model predictions into recommendations
        is deterministic and human-readable. A user (or a supervisor)
        can trace exactly why a specific recommendation was made.
\end{description}

\section{Contributions}

The concrete contributions of this project are:

\begin{itemize}
    \item A production-grade, end-to-end web application that runs three
        deep learning models in parallel and returns a report within a
        few seconds on CPU-only hardware.
    \item A comparative study of three architectural strategies --- a
        single multi-task model, three independent models with a learned
        fusion head, and three independent models with a rule-based
        engine --- that motivates our final choice.
    \item A careful treatment of class imbalance and a two-phase transfer
        learning recipe that we found to be substantially more stable on
        our small, noisy dataset than a single-phase fine-tune.
    \item A full test suite (56 tests covering unit, integration, model
        accuracy, and end-to-end pipeline) that is run before every
        change and that would have caught every regression we
        encountered during development.
    \item An honest discussion of the experiments that did not work,
        with concrete diagnoses of why they failed.
\end{itemize}

\section{Reading guide}

The book is organised so that a reader can either follow it end-to-end or
enter at whichever chapter interests them.
\Cref{ch:background} surveys the academic and industrial landscape.
\Cref{ch:academic} provides the deep-learning background needed to follow
the technical chapters.
\Cref{ch:application} describes the application from a user's point of
view.
\Cref{ch:algorithm} details our algorithm.
\Cref{ch:data} covers data and preprocessing.
\Cref{ch:experiments} presents experiments and failed experiments.
\Cref{ch:results} reports final results.
\Cref{ch:discussion,ch:conclusions,ch:future} close the book with a
discussion, conclusions, and directions for future work.

%==============================================================================
\chapter{Background and Related Work}\label{ch:background}
%==============================================================================

\section{The clinical baseline}

Modern dermatology grades acne along ordinal scales such as the
Investigator's Global Assessment (IGA) or the Global Acne Grading System
(GAGS). Both scales collapse a rich visual scene into a small number of
categories (typically \emph{clear}, \emph{mild}, \emph{moderate},
\emph{severe}). Skin type is described using variants of the classical
four-way split: \emph{oily}, \emph{dry}, \emph{normal}, \emph{combination}.
Secondary issues --- blackheads, dark spots, enlarged pores, wrinkles --- are
usually discussed qualitatively rather than graded on a formal scale.

This clinical vocabulary is important for us because it defines the target
labels our models must produce. If we invent our own categories, no
domain expert can validate our output. We therefore align our label
spaces with the categories a dermatologist would recognise.

\section{Public datasets}

Progress in dermatological deep learning has been paced by dataset
releases. The most influential public resources include:

\begin{itemize}
    \item \textbf{ACNE04}~\cite{wu2019acne04}. Approximately 1{,}457
        facial images labelled with acne severity (\emph{clear},
        \emph{mild}, \emph{moderate}, \emph{severe}) plus lesion counts.
        Images are in-the-wild and vary widely in lighting and pose,
        which is both a curse and a blessing: they are noisy but they
        resemble real user uploads.
    \item \textbf{ISIC / HAM10000}~\cite{tschandl2018ham10000}. A large
        collection of dermoscopic images labelled with skin-lesion
        categories. Excellent for melanoma / lesion classification but
        not directly applicable to cosmetic assessment because the
        images are dermoscopic rather than casual photographs.
    \item \textbf{Fitzpatrick 17k}~\cite{groh2021fitzpatrick}. A dataset
        of clinical photographs annotated with Fitzpatrick skin type,
        used mostly for fairness studies.
    \item \textbf{Kaggle Face Skin Type / Skin Problems datasets}. A
        family of smaller, community-contributed datasets covering
        skin type and skin problem categories. Quality varies but they
        are useful for filling gaps left by the more curated academic
        datasets.
\end{itemize}

We settled on ACNE04 as the anchor dataset and augmented it with
community datasets and a small in-house set of photographs. Details are
in \Cref{ch:data}.

\section{Related academic work}

Deep learning for skin analysis divides roughly into three lines of work.

\subsection{Lesion classification on dermoscopic images}
The best-known result is the work of Esteva et
al.~\cite{esteva2017dermatologist}, who trained a single Inception-v3
network on a large private dataset of clinical images and reached
dermatologist-level performance on binary melanoma classification.
Subsequent work explored specialised architectures such as
EfficientNet~\cite{tan2019efficientnet},
ResNeXt~\cite{xie2017aggregated}, and Vision Transformers, plus data
strategies such as heavy augmentation, colour constancy normalisation,
and ensemble averaging.

\subsection{Acne severity grading}
The ACNE04 paper~\cite{wu2019acne04} formulated acne severity as a
label-distribution learning problem, arguing that hard categorical
labels lose information (a photograph labelled \emph{moderate} is
often only marginally different from one labelled \emph{mild}).
Follow-up work has used ordinal regression, multi-task lesion counting,
and attention-based localisation. For our project we deliberately
chose the simpler categorical formulation, because the downstream
consumer of the prediction --- the rule engine --- expects discrete
severity buckets.

\subsection{Skin-type and cosmetic assessment}
Skin-type classification has received less academic attention than
lesion classification, and most published pipelines come from industry
rather than academia. We were unable to find a widely-cited academic
benchmark for the four-class oily/dry/normal/combination task. This is
one reason our own results on skin type are less comparable to
external numbers than our acne severity results.

\section{Commercial and industrial products}

Several commercial products have appeared in the last few years that
occupy roughly the same niche as \cara. Each one influenced our design
either by example or by counter-example.

\begin{itemize}
    \item \textbf{L'Or\'eal Skin Genius / ModiFace.} A polished mobile
        experience that returns a numerical score and product
        recommendations. The recommendations are, unsurprisingly,
        drawn exclusively from the vendor's own catalogue. This
        motivated our decision to recommend \emph{ingredients} rather
        than branded products.
    \item \textbf{Perfect Corp.\ / YouCam Makeup.} Very strong face
        tracking and augmented-reality effects. Their skin-quality
        assessment is fast but has been criticised for optimistic
        outputs that appear to be tuned for engagement rather than
        accuracy. This informed our stance on truthfulness.
    \item \textbf{Curology.} A subscription tele-dermatology service in
        which real clinicians write prescriptions after reviewing user
        photographs. Not an AI product per se, but the reference for
        what a truthful assessment looks like.
    \item \textbf{Academic demos.} A number of open-source projects on
        GitHub attempt facial skin analysis with pretrained CNNs.
        Almost all of them fall into one of two failure modes: they
        return a confident label on every input, including images with
        no face; or they degrade to noise as soon as lighting changes.
        Both failures shaped the guardrails in our pipeline.
\end{itemize}

\begin{table}[htbp]
    \centering
    \caption{Positioning of \cara relative to representative existing products.}
    \label{tab:positioning}
    \small
    \begin{tabularx}{\linewidth}{lXXXX}
        \toprule
        & \textbf{L'Or\'eal SG} & \textbf{Perfect Corp.} & \textbf{Curology} & \textbf{\cara} \\
        \midrule
        Delivery         & Mobile app       & Mobile app     & Human tele-derm  & Web app \\
        Assessment       & Learned score    & Learned score  & Human            & 3 CNNs + rules \\
        Recommends       & Own products     & Own products   & Prescription     & Ingredients \\
        Explainability   & None             & None           & Full             & Rule engine \\
        Handles unsure   & Rarely           & Rarely         & Always           & By design \\
        Open source      & No               & No             & No               & Yes \\
        \bottomrule
    \end{tabularx}
\end{table}

\section{Where \cara fits in the picture}

\cara sits deliberately in the small corner of the market that combines
(i)~open-source, reproducible, deep-learning components, with (ii)~a
transparent, rule-based recommendation layer, and (iii)~an honest
policy about uncertainty. We do not compete with L'Or\'eal on polish, and
we do not compete with Curology on medical rigour. We compete with the
countless \emph{"free online skin analysis"} demos on the accuracy and
honesty of the report they return.

%==============================================================================
\chapter{Academic Background}\label{ch:academic}
%==============================================================================

This chapter compresses the deep-learning material that the rest of the
book relies on. Readers already familiar with convolutional neural
networks, transfer learning, and standard training practice can skip
straight to \Cref{ch:application}.

\section{Image classification, formally}

An image classifier is a function
\(f_\theta : \mathcal{X} \to \Delta^{K-1}\) parameterised by weights
\(\theta\), where \(\mathcal{X}\) is the set of input images
(typically \(\mathbb{R}^{3\times H\times W}\)) and \(\Delta^{K-1}\) is
the probability simplex over \(K\) classes. Training minimises the
empirical cross-entropy loss

\begin{equation}
    \mathcal{L}(\theta) = -\frac{1}{N}\sum_{i=1}^{N} \sum_{k=1}^{K}
        w_k \, y_{i,k} \, \log f_\theta(x_i)_k
    \label{eq:xent}
\end{equation}

on a dataset \(\{(x_i, y_i)\}_{i=1}^N\), where \(y_i\) is a one-hot
target and \(w_k\) is a per-class weight that we use to counteract
class imbalance (\Cref{sec:imbalance}).

\section{Convolutional neural networks}

CNNs~\cite{lecun1998gradient} exploit two properties of images:
translation equivariance (a face is still a face when it moves five
pixels to the right) and locality (nearby pixels are correlated).
A convolutional layer computes

\begin{equation}
    Y_{c,i,j} = \sigma\!\left( b_c + \sum_{c'} \sum_{u,v}
        K_{c,c',u,v}\, X_{c',\,i+u,\,j+v} \right)
    \label{eq:conv}
\end{equation}

for each output channel \(c\), spatial position \((i,j)\) and non-linearity
\(\sigma\). Stacking convolutional layers with pooling and non-linearities
produces a hierarchy of features: edges and textures in the shallow
layers, object parts in the middle layers, and semantic categories near
the head.

\section{From AlexNet to EfficientNet}

The past decade of image classification is essentially a story of
progressively better trade-offs between depth, width, resolution, and
compute:

\begin{itemize}
    \item \textbf{AlexNet}~\cite{krizhevsky2012alexnet} (2012) proved
        that a deep CNN on GPU could crush hand-crafted features on
        ImageNet.
    \item \textbf{VGG}~\cite{simonyan2015vgg} pushed depth using small
        \(3\times 3\) filters.
    \item \textbf{ResNet}~\cite{he2016deep} introduced residual
        connections that allowed training networks with hundreds of
        layers.
    \item \textbf{Inception}~\cite{szegedy2015going} and its successors
        varied filter sizes in parallel.
    \item \textbf{EfficientNet}~\cite{tan2019efficientnet} performed a
        \emph{compound scaling} of depth, width, and input resolution
        subject to a fixed compute budget, producing a family of
        models that dominated the accuracy--FLOP frontier.
\end{itemize}

We chose the smallest member of that family, \effnet, for reasons that
are worth stating explicitly.

\begin{observation}
\effnet has approximately \num{5.3} million parameters and, at
\(224\times 224\) input, requires roughly \num{0.39} GFLOPs per
forward pass. That budget is small enough to serve three parallel
models on a CPU in under a couple of seconds, which was our
performance target.
\end{observation}

\begin{figure}[htbp]
    \centering
    \begin{tikzpicture}[
        block/.style={rectangle, draw=caraprimary!70, fill=caraprimary!8,
            minimum height=1.1cm, minimum width=1.6cm, rounded corners=2pt,
            align=center, font=\small\bfseries},
        arrow/.style={-{Latex[length=2.2mm]}, thick, draw=caraprimary!80}]
    \node[block] (in)   {Input\\224\(\times\)224};
    \node[block, right=6mm of in]    (stem) {Stem Conv};
    \node[block, right=6mm of stem]  (mb1)  {MBConv\(\times k_1\)};
    \node[block, right=6mm of mb1]   (mb2)  {MBConv\(\times k_2\)};
    \node[block, right=6mm of mb2]   (mb3)  {MBConv\(\times k_3\)};
    \node[block, below=6mm of mb2]   (gap)  {GAP};
    \node[block, right=6mm of gap]   (fc)   {Dropout + FC};
    \node[block, right=6mm of fc]    (out)  {Softmax\\\(K\) classes};

    \draw[arrow] (in) -- (stem);
    \draw[arrow] (stem) -- (mb1);
    \draw[arrow] (mb1) -- (mb2);
    \draw[arrow] (mb2) -- (mb3);
    \draw[arrow] (mb3.south) |- (gap.east);
    \draw[arrow] (gap) -- (fc);
    \draw[arrow] (fc) -- (out);
    \end{tikzpicture}
    \caption{Schematic of the EfficientNet-B0 backbone used for each
    of \cara's three classifiers. The classifier head (dropout +
    fully-connected + softmax) is the only part re-initialised for
    each task.}
    \label{fig:effnet-block}
\end{figure}

\section{Transfer learning}

Training a modern CNN from random weights on 1{,}000 images is a
recipe for overfitting. Transfer learning
\cite{yosinski2014transferable} solves this by initialising the network
with weights pretrained on a large source dataset --- for us,
ImageNet-1k~\cite{deng2009imagenet} --- and then adapting the model to
the target task. Two knobs control the process:

\begin{itemize}
    \item Which layers to freeze and for how long.
    \item How the learning rate differs between the pretrained backbone
        and the newly-initialised head.
\end{itemize}

Our two-phase recipe (\Cref{sec:two-phase}) is a straightforward
instance of this pattern.

\section{Regularisation and augmentation}

Small datasets punish confident networks. We rely on three
regularisers:

\begin{enumerate}
    \item \textbf{Data augmentation.} Random horizontal flips, mild
        rotations, colour jitter, and random cropping. All
        augmentations are applied to training data only, never to
        validation or test data.
    \item \textbf{Dropout} in the classifier head.
    \item \textbf{Weight decay} via AdamW~\cite{loshchilov2019adamw}.
\end{enumerate}

\section{Class imbalance}\label{sec:imbalance}

Real-world skin datasets are heavily skewed. In ACNE04, our
\emph{severe} class contains roughly one seventh as many samples as
\emph{mild}. Naive training would produce a classifier that predicts
\emph{mild} everywhere and reports respectable accuracy.

We use weighted cross-entropy with weights derived from inverse class
frequency:

\begin{equation}
    w_k = \frac{N}{K \cdot n_k}
    \label{eq:weights}
\end{equation}

where \(N\) is the total number of training samples, \(K\) the number
of classes, and \(n_k\) the number of samples in class \(k\). The
weights we actually used are listed in \Cref{ch:data}.

\section{Optimisation}

We use AdamW with an initial learning rate of \(10^{-3}\) for the
classifier head and \(10^{-4}\) for the unfrozen backbone. Learning
rate scheduling is cosine annealing with a short linear warm-up.
Early stopping is triggered by validation loss, not validation
accuracy, because loss is more sensitive to the confidence calibration
we care about.

\section{Rule-based systems and BLP}

The last piece of academic background is on the recommendation side.
Business Logic Programming (BLP) is our internal name for the
data-driven, deterministic rule engine that converts model outputs
into user-facing advice. It is closer in spirit to a decision table
than to a formal rule system such as Datalog, but the important point
is that every recommendation is produced by evaluating a hand-written
rule against the model's structured prediction. There is no
learning component in the recommendation step, by design: any
recommendation shown to the user can be justified by pointing at a
single rule in a JSON file.

\section{Backpropagation, in one page}

The training loop of \cara is the same textbook backpropagation loop
used by every modern deep learning framework, and we assume the
reader has met it before. For completeness we summarise it here.

Given a loss \(\mathcal{L}(\theta)\) that is a scalar function of the
model parameters \(\theta\), backpropagation computes
\(\nabla_\theta \mathcal{L}\) by repeated application of the chain
rule. If the network is a composition
\(f = f_L \circ f_{L-1} \circ \cdots \circ f_1\) with intermediate
activations \(a_\ell = f_\ell(a_{\ell-1}; \theta_\ell)\), then

\begin{equation}
    \frac{\partial \mathcal{L}}{\partial \theta_\ell}
    \;=\;
    \frac{\partial \mathcal{L}}{\partial a_L}
    \cdot
    \prod_{k=\ell+1}^{L} \frac{\partial a_k}{\partial a_{k-1}}
    \cdot
    \frac{\partial a_\ell}{\partial \theta_\ell}.
    \label{eq:backprop}
\end{equation}

The parameter update follows some variant of stochastic gradient
descent. We use AdamW, whose update rule adds momentum and
per-parameter adaptive scaling:

\begin{align}
    m_t &= \beta_1 m_{t-1} + (1-\beta_1) g_t, \\
    v_t &= \beta_2 v_{t-1} + (1-\beta_2) g_t^2, \\
    \hat{m}_t &= m_t / (1-\beta_1^t), \\
    \hat{v}_t &= v_t / (1-\beta_2^t), \\
    \theta_{t+1} &= \theta_t - \eta \left(
        \frac{\hat{m}_t}{\sqrt{\hat{v}_t} + \epsilon} + \lambda \theta_t
    \right).
    \label{eq:adamw}
\end{align}

The decoupled weight-decay term \(\lambda \theta_t\) is what
distinguishes AdamW from vanilla Adam and is a significant part of
why it generalises better in our experiments.

\section{Softmax, cross-entropy and their gradient}

Every one of our classifiers ends with a softmax layer

\begin{equation}
    p_k = \frac{\exp(z_k)}{\sum_{j=1}^{K}\exp(z_j)},
    \label{eq:softmax}
\end{equation}

whose input \(z\) is the logit vector produced by the final linear
layer. Cross-entropy loss against a one-hot target \(y\) evaluates to
\(\mathcal{L} = -\log p_{k^*}\) where \(k^*\) is the index of the
true class. A short calculation shows that its gradient with respect
to logits is the intuitive expression \(\partial \mathcal{L}/\partial
z_k = p_k - y_k\), which is the reason the softmax--cross-entropy
combination is numerically and computationally so pleasant.

\section{Batch normalisation and its role in EfficientNet}

Every stage of an EfficientNet contains batch normalisation layers
\cite{ioffe2015batchnorm}, which whiten the intermediate activations
before the non-linearity:

\begin{equation}
    \hat{a}_{c} = \gamma_c \cdot \frac{a_c - \mu_c}{\sqrt{\sigma_c^2 + \epsilon}} + \beta_c.
    \label{eq:bn}
\end{equation}

Two consequences matter for our project. First, batch norm reduces
the sensitivity to the initial learning rate, which stabilised
fine-tuning in our two-phase recipe. Second, the running mean and
variance are estimated on the fly during training and then frozen at
inference time; forgetting to switch the model to \texttt{eval()}
mode before inference was, in fact, one of our early bugs and cost us
a full day of debugging before we spotted the regression on the test
set.

\section{Activation functions}

\effnet uses the Swish (a.k.a.\ SiLU) activation
\(\sigma(x) = x \cdot \operatorname{sigmoid}(x)\), which is smoother
than ReLU and empirically slightly more accurate. We did not change
this choice. In the classifier head we deliberately used ReLU because
the head is small and the interpretability of ReLU is a small
convenience during debugging.

\section{Learning-rate scheduling}

We use a cosine annealing schedule \cite{loshchilov2017sgdr} with a
short linear warm-up:

\begin{equation}
    \eta_t = \eta_{\min} + \tfrac{1}{2}(\eta_{\max} - \eta_{\min})
    \left( 1 + \cos\!\left( \pi \cdot t / T \right) \right).
    \label{eq:cosine}
\end{equation}

Cosine annealing is a small choice with an outsized effect on the
final validation-accuracy plateau. Constant learning rates in early
experiments left the model oscillating around a suboptimal minimum;
cosine-annealed runs converged to a noticeably lower validation
loss and a smoother loss curve.

\section{Metrics beyond accuracy}

We report accuracy as the headline metric because it is what a
non-expert user expects, but internally we track a small dashboard of
secondary metrics that we consider more informative on imbalanced
data:

\begin{itemize}
    \item \textbf{Macro-averaged \(F_1\).} The unweighted mean of the
        per-class \(F_1\) scores. Insensitive to class imbalance.
    \item \textbf{Balanced accuracy.} The unweighted mean of per-class
        recall. Comparable to accuracy on balanced data but more
        honest on imbalanced data.
    \item \textbf{Confusion matrix.} The full per-class breakdown.
        Any conversation about where a model is failing starts here.
    \item \textbf{Expected calibration error (ECE)}. Averaged
        difference between confidence and accuracy across bins,
        introduced by Guo et al.~\cite{guo2017calibration}.
\end{itemize}

We look at ECE alongside accuracy for the same reason we introduced
the confidence-based ``Uncertain'' rendering in the frontend: an
overconfident wrong answer is worse than a hesitant right answer.

%==============================================================================
\chapter{The Application}\label{ch:application}
%==============================================================================

\section{User story}

The primary user of \cara is a curious, non-expert consumer who has a
skincare question they cannot easily answer. A typical interaction
takes about a minute.

\begin{enumerate}
    \item The user opens the web app on a phone or a laptop and clicks
        \emph{Analyze}.
    \item They upload a well-lit selfie or take one directly with the
        camera.
    \item The frontend uploads the image to the backend over a POST
        request.
    \item Within a few seconds, the backend responds with a structured
        report.
    \item The frontend renders the report: acne severity, skin type,
        detected issues, confidence bars, ingredients to look for,
        ingredients to avoid, and a plain-English explanation.
    \item The report is saved to storage. The user can browse past
        reports on a \emph{History} page.
\end{enumerate}

\section{System architecture}

\Cref{fig:architecture} shows the runtime architecture. Everything
inside the dashed box runs inside a single FastAPI process, which we
package into a Docker container for deployment. The React frontend is
served from an nginx container. In production, both containers sit
behind an nginx reverse proxy.

\begin{figure}[htbp]
    \centering
    \begin{tikzpicture}[
        node distance=6mm and 10mm,
        every node/.style={font=\small},
        service/.style={rectangle, draw=caraprimary, fill=caraprimary!8,
            minimum height=1.0cm, minimum width=3.2cm, rounded corners=2pt,
            align=center},
        model/.style={rectangle, draw=caraaccent, fill=caraaccent!10,
            minimum height=0.9cm, minimum width=3.2cm, rounded corners=2pt,
            align=center},
        io/.style={ellipse, draw=black!60, fill=black!5, align=center,
            minimum height=0.9cm, minimum width=2.6cm},
        arrow/.style={-{Latex[length=2.5mm]}, thick, draw=black!70}]

        \node[io]                        (user)  {User browser};
        \node[service, below=of user]    (front) {React SPA\\(Home / Analyze / Results / History)};
        \node[service, below=of front]   (api)   {FastAPI Routes};
        \node[service, below=of api]     (svc)   {Analysis Service};
        \node[service, below=of svc]     (pre)   {Preprocessing pipeline\\(face detect \(\to\) CLAHE \(\to\) resize)};

        \node[service, below=8mm of pre] (orch)  {Inference Orchestrator};
        \node[model, below left=6mm and -10mm of orch] (m1) {Acne severity\\EfficientNet-B0};
        \node[model, below=6mm of orch]                (m2) {Skin type\\EfficientNet-B0};
        \node[model, below right=6mm and -10mm of orch](m3) {Skin issues\\EfficientNet-B0};

        \node[service, below=22mm of m2] (blp)  {BLP Rule Engine\\(rules.json)};
        \node[service, below=of blp]     (rep)  {Report Builder\\+ Storage};

        \draw[arrow] (user)  -- (front);
        \draw[arrow] (front) -- (api);
        \draw[arrow] (api)   -- (svc);
        \draw[arrow] (svc)   -- (pre);
        \draw[arrow] (pre)   -- (orch);
        \draw[arrow] (orch)  -- (m1);
        \draw[arrow] (orch)  -- (m2);
        \draw[arrow] (orch)  -- (m3);
        \draw[arrow] (m1) |- (blp);
        \draw[arrow] (m2) -- (blp);
        \draw[arrow] (m3) |- (blp);
        \draw[arrow] (blp)   -- (rep);
        \draw[arrow, dashed] (rep.east) .. controls +(right:2.5cm) and +(right:2.5cm) .. (front.east);

        \begin{scope}[on background layer]
            \node[draw=black!40, dashed, rounded corners=3pt,
                  fit=(api)(svc)(pre)(orch)(m1)(m2)(m3)(blp)(rep),
                  inner sep=6pt, label={[font=\footnotesize\itshape]above:Backend (FastAPI + PyTorch)}] {};
        \end{scope}
    \end{tikzpicture}
    \caption{End-to-end runtime architecture of \cara.}
    \label{fig:architecture}
\end{figure}

\section{Frontend}

The frontend is a React 18 single-page application. There are four
pages: \emph{Home} (marketing landing page and disclaimer),
\emph{Analyze} (image upload and camera capture), \emph{Results} (the
report), and \emph{History} (list of past reports). Styling is
mobile-first because most casual selfies come from phones.

We deliberately kept the frontend thin. It knows how to upload an
image, poll for results, and render JSON. All logic --- including
confidence bands and colour coding --- is driven by fields returned in
the response payload. This makes it possible to iterate on the model
side without redeploying the frontend.

\section{Backend}

The backend is a FastAPI application. FastAPI was chosen for three
reasons: automatic OpenAPI schema generation (which makes the frontend
integration typed and testable), Pydantic-based request validation
(which rejects malformed inputs cheaply), and native async support
(which we do not use aggressively yet but which leaves room for
future concurrency).

The main endpoint is:

\begin{lstlisting}[language=Python,caption={The main analysis endpoint (simplified).}]
@router.post("/analyze", response_model=AnalysisReport)
async def analyze(image: UploadFile = File(...)) -> AnalysisReport:
    image_bytes = await image.read()
    report = await analysis_service.analyze(image_bytes)
    return report
\end{lstlisting}

The \texttt{AnalysisService} orchestrates the pipeline: preprocessing,
inference, confidence checking, BLP evaluation, report building, and
persistence. Each step is a separate injectable component so that we
can mock any of them during testing.

\section{Storage}

Reports are persisted as JSON files under a \texttt{storage/}
directory. This is a deliberately boring choice: it keeps the demo
runnable without any external database and it makes reports easy to
inspect during debugging. Behind an abstract \texttt{StorageRepository}
interface, we ship a second implementation for MongoDB (via Motor)
that can be switched on with an environment variable. We validated
the Mongo path in a dev container but the deployed instance uses JSON.

\section{Deployment}

The project ships with a \texttt{docker-compose.yml} that starts three
containers: the FastAPI backend, the nginx-fronted React app, and a
tiny nginx reverse proxy. Model weights are baked into the backend
image at build time so that a container start does not depend on
network access.

\section{Error handling and user experience}

Each layer of the stack has a defined failure policy.

\begin{itemize}
    \item \textbf{Preprocessing.} If no face is detected the pipeline
        returns a structured error that the frontend renders as
        ``No face detected. Try better lighting or a closer shot.'' We
        never fall back to running the model on the whole image.
    \item \textbf{Inference.} Every model prediction includes a
        confidence score. If any model is below a per-model
        confidence floor, the corresponding section of the report is
        labelled ``Uncertain'' rather than being replaced by the
        top-1 class.
    \item \textbf{API.} The endpoint returns typed error responses
        with error codes; raw tracebacks are never exposed to the
        user.
    \item \textbf{Frontend.} User-facing messages are actionable,
        never technical.
\end{itemize}

\section{Frontend, in more depth}

The React frontend is organised around a small state machine. The
possible states are:

\begin{center}
\begin{tikzpicture}[
    node distance=8mm and 12mm,
    every node/.style={font=\footnotesize},
    state/.style={rectangle, draw=caraprimary, fill=caraprimary!8,
        rounded corners=2pt, minimum height=0.9cm, minimum width=2.2cm,
        align=center},
    err/.style={rectangle, draw=caraaccent, fill=caraaccent!10,
        rounded corners=2pt, minimum height=0.9cm, minimum width=2.2cm,
        align=center},
    arrow/.style={-{Latex[length=2mm]}, thick, draw=black!70}]
\node[state]                    (idle)     {Idle};
\node[state, right=of idle]     (selected) {Image\\selected};
\node[state, right=of selected] (upload)   {Uploading};
\node[state, right=of upload]   (waiting)  {Waiting for\\report};
\node[state, right=of waiting]  (done)     {Report\\shown};
\node[err, below=of upload]     (error)    {Error};
\draw[arrow] (idle)     -- (selected);
\draw[arrow] (selected) -- (upload);
\draw[arrow] (upload)   -- (waiting);
\draw[arrow] (waiting)  -- (done);
\draw[arrow] (done.south)   .. controls +(down:8mm) and +(down:8mm) .. (idle.south);
\draw[arrow] (upload.south)  -- (error.north);
\draw[arrow] (error.west)  .. controls +(left:8mm) and +(down:8mm)  .. (idle.south);
\end{tikzpicture}
\end{center}

Only one screen is rendered at a time; navigation is handled by
React Router. The state machine is deliberately simple because most
of the interesting logic lives on the backend.

The report page renders each section as a card. Each card shows the
predicted label, a colour-coded confidence bar, and a short
plain-English explanation string returned by the backend. If the
backend flagged a section as \emph{uncertain}, the card displays a
neutral colour and the label ``Uncertain --- try a better-lit
photograph''. This is the only place in the frontend where user
copy is not entirely driven by the API response, and we intend to
push it into the backend as part of future work.

\section{Backend, in more depth}

The backend is organised around three layers.

\begin{description}
    \item[Routes.] Thin HTTP handlers under \texttt{backend/api/}.
        They validate the request via Pydantic, hand off to a
        service, and marshal the response. Routes contain no
        business logic.
    \item[Services.] The main entry point is
        \texttt{AnalysisService}, which orchestrates a single
        request. \texttt{ReportBuilder} handles the assembly of the
        final \texttt{AnalysisReport} object.
    \item[Repositories and engines.] The
        \texttt{InferenceOrchestrator}, the BLP engine, and the
        \texttt{StorageRepository} are injectable dependencies. Each
        has a well-defined interface and each is mocked in tests.
\end{description}

Dependency injection is done with FastAPI's built-in
\texttt{Depends} mechanism. It is enough for our size; we resisted
the urge to introduce a heavier DI framework.

%==============================================================================
\chapter{The Algorithm}\label{ch:algorithm}
%==============================================================================

\section{Overview}

The core algorithm of \cara is a three-stage pipeline: preprocessing,
inference, and rule-based recommendation. \Cref{fig:algorithm} shows
the data flow. This chapter walks through each stage in detail.

\begin{figure}[htbp]
    \centering
    \begin{tikzpicture}[
        node distance=10mm,
        every node/.style={font=\small},
        stage/.style={rectangle, draw=caraprimary!70, fill=caraprimary!8,
            minimum height=1.6cm, minimum width=3.4cm, rounded corners=2pt,
            align=center},
        data/.style={rectangle, draw=black!60, fill=black!5,
            minimum height=1.0cm, minimum width=2.8cm, align=center},
        arrow/.style={-{Latex[length=2.5mm]}, thick, draw=black!70}]

        \node[data]                        (raw)   {Raw image};
        \node[stage, right=of raw]         (pre)   {Preprocess\\face crop, CLAHE, resize};
        \node[data,  right=of pre]         (tens)  {Tensor \(3\times224\times224\)};
        \node[stage, below=of tens]        (inf)   {Inference\\3 EfficientNets};
        \node[data,  left=of inf]          (preds) {Structured predictions};
        \node[stage, left=of preds]        (blp)   {BLP engine\\rules.json};
        \node[data,  below=of blp]         (rep)   {Report};

        \draw[arrow] (raw)   -- (pre);
        \draw[arrow] (pre)   -- (tens);
        \draw[arrow] (tens)  -- (inf);
        \draw[arrow] (inf)   -- (preds);
        \draw[arrow] (preds) -- (blp);
        \draw[arrow] (blp)   -- (rep);
    \end{tikzpicture}
    \caption{High-level data flow of the \cara algorithm.}
    \label{fig:algorithm}
\end{figure}

\section{Preprocessing}\label{sec:preprocessing}

Preprocessing is where a large fraction of our engineering effort went,
and probably the largest single contributor to accuracy after
architecture choice.

\subsection{Face detection and cropping}
Uploaded images vary wildly. Some are tight selfies; some are group
photos; some are landscape shots with a face in the corner. Feeding
the raw image to the model would waste most of the receptive field on
background pixels. We therefore detect the face first and crop with a
20\% padding margin before resizing.

We evaluated three detectors: OpenCV Haar cascade, MediaPipe Face
Detection, and MTCNN. The trade-offs (\Cref{tab:face-detect}) drove
our decision to use Haar as the default with an explicit fallback
plan.

\begin{table}[htbp]
    \centering
    \caption{Face detector trade-offs on our internal test set of 200 selfies.}
    \label{tab:face-detect}
    \small
    \begin{tabularx}{\linewidth}{lXXX}
        \toprule
        \textbf{Detector} & \textbf{Recall} & \textbf{Speed (CPU)} & \textbf{Comment} \\
        \midrule
        Haar cascade   & Moderate  & Very fast    & Fails on side profiles; sufficient for well-framed selfies. \\
        MediaPipe      & High      & Fast         & Adds a moderate dependency; behaves well across lighting. \\
        MTCNN          & Very high & Slow on CPU  & Overkill for our use case; noticeable latency. \\
        \bottomrule
    \end{tabularx}
\end{table}

\subsection{Brightness normalisation with CLAHE}
Even after cropping, lighting varies enormously. Naive histogram
equalisation destroys skin colour and confuses the skin-type model.
We use CLAHE (Contrast-Limited Adaptive Histogram Equalisation) on
the L channel of the LAB colour space, which normalises local
brightness without destroying colour information.

\subsection{Resizing and tensor normalisation}
Finally we resize the cropped face to \(224\times 224\), convert to a
\pytorch tensor, and normalise with the standard ImageNet mean and
standard deviation. This last step is important because our
transfer-learning weights were trained on ImageNet-normalised inputs.

\section{Inference}

\subsection{Why three models and not one}
Our first prototype used a single \effnet with a multi-head output
covering all three tasks. It was clean, fast, and worse than three
independent models on every task. We diagnose the failure in
\Cref{sec:failed-monolithic}. The short version: the three tasks
compete for capacity in a small backbone, and the loss weights that
balance them are difficult to tune without validation for each task.

Our final design is three independent \effnet models, each fine-tuned
on its own dataset with its own head. They share the same
preprocessing pipeline and the same input tensor.

\subsection{The Model Registry pattern}
Loading model weights is the slowest part of a cold start, so we do it
once. A \texttt{ModelRegistry} singleton is initialised at FastAPI
startup and holds the three models in memory. The
\texttt{InferenceOrchestrator} pulls models from the registry and runs
predictions; it never touches disk during a request.

\begin{lstlisting}[language=Python,caption={Sketch of the InferenceOrchestrator.}]
class InferenceOrchestrator:
    def __init__(self, registry: ModelRegistry):
        self._registry = registry

    def predict_all(self, tensor: torch.Tensor) -> AllPredictions:
        return AllPredictions(
            acne       = self._predict(self._registry.acne, tensor),
            skin_type  = self._predict(self._registry.skin_type, tensor),
            skin_issue = self._predict(self._registry.skin_issue, tensor),
        )

    def _predict(self, model, tensor) -> ModelPrediction:
        with torch.no_grad():
            logits = model(tensor)
            probs  = torch.softmax(logits, dim=-1)
            conf, idx = probs.max(dim=-1)
        return ModelPrediction(
            label      = model.classes[idx.item()],
            confidence = conf.item(),
            probs      = probs.squeeze().tolist(),
        )
\end{lstlisting}

\subsection{Confidence and uncertainty}
Every prediction carries a confidence score (the maximum softmax
probability). Downstream code treats confidence as a first-class
signal: predictions below a per-model floor are rendered as
``Uncertain'' rather than as a hard label. This is the mechanism that
lets us keep our promise of truthfulness.

\section{The BLP recommendation engine}

The BLP engine takes a structured prediction and evaluates a set of
rules loaded from a JSON file. Each rule has a precondition
(a partial specification of the prediction such as ``\texttt{acne:
moderate} and \texttt{skin\_type: oily}''), a set of ingredient
recommendations to promote, a set of ingredients to warn against, and
a plain-English explanation string.

\begin{lstlisting}[language=Python,caption={A representative BLP rule.}]
{
  "id": "moderate_acne_oily_skin",
  "when": {
    "acne_severity": ["moderate", "severe"],
    "skin_type": ["oily", "combination"]
  },
  "recommend_ingredients": [
    "salicylic acid (BHA)",
    "niacinamide",
    "benzoyl peroxide (spot use)"
  ],
  "avoid_ingredients": [
    "heavy occlusives",
    "coconut oil"
  ],
  "explanation": "Excess sebum and active acne benefit from a leave-on BHA to keep pores clear and niacinamide to reduce inflammation. Heavy occlusives trap sebum."
}
\end{lstlisting}

Rules are matched in a deterministic order. The engine returns the
union of all matching rules, deduplicated. There is no learning
component --- we can prove that the recommendation set for a given
prediction is a pure function of the prediction and the rules file.

\subsection{Anatomy of a rule}

Each rule has five fields.

\begin{description}
    \item[\texttt{id}] A stable, human-readable identifier. It appears
        in the report metadata so that a future maintainer can
        trace a specific recommendation back to the specific rule
        that produced it.
    \item[\texttt{when}] A partial specification of the prediction
        that the rule matches. Fields not listed in \texttt{when}
        are wildcards. This is the mechanism that lets a single
        rule handle a family of related predictions.
    \item[\texttt{recommend\_ingredients}] Ingredients to look for.
        These become the ``Look for'' section of the report.
    \item[\texttt{avoid\_ingredients}] Ingredients to avoid. These
        become the ``Avoid'' section of the report.
    \item[\texttt{explanation}] A short plain-English string that is
        surfaced directly in the report. Explanations are written
        for a non-expert reader.
\end{description}

\subsection{Rule design principles}

We authored the rules under three self-imposed constraints.

\begin{enumerate}
    \item \textbf{Every rule is defensible.} Each ingredient
        recommendation is supported by consumer-facing dermatology
        literature (peer-reviewed sources where possible, reputable
        dermatology textbooks otherwise). We keep a companion
        bibliography of the sources we consulted, one per rule.
    \item \textbf{No conflicting recommendations.} If two rules
        match the same prediction, they may not recommend an
        ingredient from one rule while the other rule warns against
        it. The BLP test suite includes a lint that fails the build
        if this invariant is ever violated.
    \item \textbf{Uncertainty short-circuits recommendations.} When
        one of the models is flagged as uncertain, the corresponding
        section of the report is filled with a neutral message
        rather than a rule-based recommendation. We would rather
        say ``we could not confidently classify this'' than
        recommend niacinamide to somebody whose skin type we do not
        actually know.
\end{enumerate}

\subsection{Coverage matrix}

The rules file must cover the Cartesian product of the three model
label spaces, or at least, every subset that is likely to appear.
\Cref{tab:coverage} sketches how the rules are distributed across
severity by skin type combinations for the acne / skin-type
sub-space. Zero-cell entries fall back to a default rule that
recommends a gentle cleanser and moisturiser and warns against
harsh actives.

\begin{table}[htbp]
    \centering
    \caption{Number of BLP rules matching each (acne, skin type)
    combination in the current rules file.}
    \label{tab:coverage}
    \small
    \begin{tabular}{lcccc}
        \toprule
        & \textbf{oily} & \textbf{dry} & \textbf{normal} & \textbf{combination} \\
        \midrule
        clear    & 1 & 1 & 1 & 1 \\
        mild     & 2 & 2 & 1 & 2 \\
        moderate & 3 & 2 & 2 & 3 \\
        severe   & 3 & 2 & 2 & 3 \\
        \bottomrule
    \end{tabular}
\end{table}

\section{Two-phase transfer learning}\label{sec:two-phase}

Each of the three models is trained with the following two-phase
recipe.

\begin{description}
    \item[Phase 1 --- head only.] The entire backbone is frozen. Only
        the newly-initialised classifier head is trained, with a
        moderate learning rate (typically \num{1e-3}). This lets the
        head learn to project ImageNet features onto the new label
        space without disturbing them.
    \item[Phase 2 --- full fine-tune.] The backbone is unfrozen. The
        head keeps its learning rate; the backbone gets a smaller
        learning rate (typically \num{1e-4}). This lets the model
        adapt low-level features to the specifics of skin images.
\end{description}

Empirically, phase 1 alone gets us most of the way to the final
accuracy; phase 2 typically adds another 6--10 percentage points and
also sharpens the confidence calibration.

\Cref{fig:val-curves} shows a representative validation curve. Both
phases are visible: the first ends when the frozen-backbone accuracy
plateaus, and the second continues from where the first stopped.

\begin{figure}[htbp]
    \centering
    \begin{tikzpicture}
    \begin{axis}[
        width=\linewidth, height=6.5cm,
        xlabel={Epoch}, ylabel={Validation accuracy},
        ymin=0.3, ymax=0.85,
        xmin=0, xmax=25,
        grid=both, grid style={gray!20},
        legend pos=south east,
        legend style={font=\footnotesize},
        title={Two-phase training on the acne severity task}]
    \addplot[thick, color=caraprimary, mark=*, mark size=1.4pt] coordinates {
        (1,0.469) (2,0.551) (3,0.614) (4,0.628) (5,0.652)
    };
    \addlegendentry{Phase 1 (head only)}
    \addplot[thick, color=caraaccent, mark=square*, mark size=1.4pt] coordinates {
        (6,0.647) (7,0.686) (8,0.696) (9,0.715) (10,0.729)
    };
    \addlegendentry{Phase 2 (full fine-tune)}
    \addplot[dashed, color=black!60] coordinates {(0,0.696) (25,0.696)};
    \addlegendentry{Final test accuracy}
    \end{axis}
    \end{tikzpicture}
    \caption{Representative validation-accuracy curve on the acne
    severity task. Numbers taken from our training log.}
    \label{fig:val-curves}
\end{figure}

%==============================================================================
\chapter{Software Engineering, Testing, and MLOps}\label{ch:engineering}
%==============================================================================

A deep-learning project book usually devotes ninety per cent of its
attention to the models. We resist that convention because, in the
day-to-day life of \cara, the engineering scaffolding around the
models mattered as much as the models themselves. This chapter
documents the choices that let three students on a laptop iterate
without breaking each other's work.

\section{Repository layout}

The repository is organised as a monorepo containing four cooperating
subsystems:

\begin{itemize}
    \item \texttt{backend/} --- the FastAPI application, including
        the analysis service, BLP engine, storage repository and
        report builder.
    \item \texttt{model\_service/} --- everything related to the deep
        learning models: architecture definitions, preprocessing
        pipeline, model registry, training scripts, and dataset
        loaders. \texttt{checkpoints/} is git-ignored to keep the
        repository small.
    \item \texttt{frontend/} --- the React SPA and its build
        artifacts.
    \item \texttt{shared/} --- Pydantic schemas and configuration
        that is imported from both the backend and the model
        service. Keeping the schemas in a single place prevents the
        canonical drift between frontend expectations and backend
        responses.
\end{itemize}

Tests live in a top-level \texttt{tests/} folder rather than being
scattered next to the code they exercise. This is a deliberate
choice: it makes it easy to run the full suite as a gate, and it
discourages the anti-pattern of ``I'll write the test later''.

\section{Configuration and secrets}

There are no secrets in this project. Model weights are shipped with
the container image; the JSON storage backend needs no credentials;
the MongoDB backend, when enabled, reads its connection string from
an environment variable that is documented in \texttt{.env.example}.
This is deliberate: a small team is more likely to leak credentials
than to lose data, and every configuration knob we add is a
configuration knob a future maintainer must understand.

\section{Testing pyramid}

The test suite is organised as a classical pyramid.

\begin{table}[htbp]
    \centering
    \caption{Composition of the \cara test suite (56 tests total).}
    \label{tab:tests}
    \small
    \begin{tabular}{lrp{7cm}}
        \toprule
        \textbf{File} & \textbf{Tests} & \textbf{What it verifies} \\
        \midrule
        \texttt{test\_blp.py}             & 18 & BLP rule engine logic, that every rule is reachable, no rule ties. \\
        \texttt{test\_api.py}             &  4 & API endpoints, request/response schema, error paths. \\
        \texttt{test\_integration.py}     &  8 & Component wiring, report assembly, storage round-trip. \\
        \texttt{test\_model\_accuracy.py} & 11 & Per-model accuracy on curated ground-truth images. \\
        \texttt{test\_e2e\_pipeline.py}   & 10 & Full image-to-report pipeline on real images. \\
        \midrule
        \textbf{Total} & \textbf{56} & \\
        \bottomrule
    \end{tabular}
\end{table}

\subsection{Unit tests}
The 18 BLP tests are the fastest and the strictest. They lock down
the behaviour of the rule engine on a matrix of representative
inputs. Every time we edit the rules file, we also touch these
tests. If a rule change unintentionally silences a recommendation,
we find out within seconds.

\subsection{Integration tests}
Integration tests stub out the models with deterministic fake
predictors and exercise the report builder, the storage layer, and
the API contract. They run in a fraction of a second and catch
schema drift bugs before they reach the frontend.

\subsection{Model accuracy tests}
The 11 model accuracy tests are the most expensive and the most
valuable. Each test loads a specific labelled image from a small
ground-truth folder and asserts that the corresponding model
predicts the correct top-1 label. Together they form a small but
brutally honest regression suite: a subtle bug in preprocessing that
degrades accuracy by one image will fail one of these tests.

\subsection{End-to-end tests}
The 10 end-to-end tests exercise the full pipeline: real photograph
in, real report out. They also assert on the shape of the
recommendation set, not just the model output, so they double as
regression tests for the BLP engine.

\section{Continuous integration}

The GitHub Actions workflow runs \texttt{pytest} on every push and
every pull request. If the model accuracy suite fails, the merge is
blocked. This is by policy: a green build is a green pipeline, not a
green build with a comment saying ``model is a bit worse but we'll
fix it later''.

\section{Reproducibility}

We fix random seeds for Python, NumPy, and \pytorch to \num{42}. We
pin dependency versions in \texttt{requirements.txt}. We pin the
CUDA-enabled variant of PyTorch separately in the training
documentation. On the same hardware, our headline accuracy numbers
are reproducible to within a few thousandths.

\section{Debugging techniques we relied on}

Three debugging patterns paid for themselves many times over.

\begin{enumerate}
    \item \textbf{Tensor shape assertions everywhere.} Every function
        that receives a tensor asserts its expected shape on entry.
        Silent shape bugs are the most vicious class of deep-learning
        bug; a two-line assertion turns them into loud failures.
    \item \textbf{Deterministic fake models in tests.} The
        \texttt{ModelRegistry} accepts a fake model in tests. That
        lets us run the entire pipeline without the CNNs, which
        speeds up the loop and pinpoints regressions in the plumbing
        rather than the models.
    \item \textbf{Storage as a debugging surface.} Because every
        report is a JSON file, reproducing a user report is a
        one-liner: load the JSON, feed the image back through the
        pipeline, compare. We used this to diagnose more than one
        regression in the wild.
\end{enumerate}

\section{Deployment story}

\cara is packaged as three Docker containers behind an nginx reverse
proxy. Model weights are baked into the backend image at build time,
which trades a slightly larger image for cold-start reliability.
Health-check endpoints let Docker Compose restart a stuck backend
without human intervention. There is no autoscaling and no cloud
dependency; the whole system runs happily on a modest VM.

\section{What we would do differently next time}

If we started the project again on day one, we would:

\begin{itemize}
    \item Freeze the label spaces \emph{before} downloading any
        data. Every dataset we pulled came with slightly different
        label names, and normalising them was a surprising amount of
        work.
    \item Write the test suite alongside the first prototype rather
        than after it. Retrofitting tests onto working code is
        harder than writing them first.
    \item Adopt Weights \& Biases (or a similar experiment tracker)
        from day one. Our home-grown training log is readable but
        cumbersome to filter.
    \item Use MediaPipe from the beginning rather than Haar. We got
        away with Haar for well-framed selfies, but real user
        uploads exercise more of the corner cases MediaPipe handles
        cleanly.
\end{itemize}

%==============================================================================
\chapter{Data and Preprocessing}\label{ch:data}
%==============================================================================

\section{Dataset composition}

Our training data is a union of three sources.

\begin{itemize}
    \item \textbf{ACNE04}~\cite{wu2019acne04}, our anchor dataset for
        acne severity.
    \item Community datasets from Kaggle for skin type and skin
        issues.
    \item A small in-house set of photographs, contributed by team
        members and volunteers, used mainly for the acceptance and
        end-to-end tests.
\end{itemize}

\section{Splits}

We use a 70/15/15 train / validation / test split, generated by a
stratified random sampler so that class proportions are preserved in
each split. Splits are fixed once and stored on disk; we never
re-generate them during training. This is a small detail that
matters: it means that a validation curve from week~3 is directly
comparable to one from week~8.

\begin{table}[htbp]
    \centering
    \caption{Split sizes for the acne severity task (ACNE04-based).}
    \label{tab:splits-acne}
    \small
    \begin{tabular}{lrrrr}
        \toprule
        & \textbf{Total} & \textbf{Train} & \textbf{Val} & \textbf{Test} \\
        \midrule
        Images & 1\,377 & 963 & 207 & 207 \\
        \bottomrule
    \end{tabular}
\end{table}

\begin{table}[htbp]
    \centering
    \caption{Class distribution and derived per-class weights for the
    acne severity training set.}
    \label{tab:acne-weights}
    \small
    \begin{tabular}{lrrr}
        \toprule
        \textbf{Class} & \textbf{Count} & \textbf{Frequency} & \textbf{Weight \(w_k\)} \\
        \midrule
        clear    & 338 & 35.1\% & 0.712 \\
        mild     & 436 & 45.3\% & 0.552 \\
        moderate & 122 & 12.7\% & 1.973 \\
        severe   &  67 &  7.0\% & 3.593 \\
        \midrule
        \textbf{Total} & 963 & 100.0\% & --- \\
        \bottomrule
    \end{tabular}
\end{table}

\section{Cleaning}

Community datasets are noisy. We spent a substantial amount of time
in manual triage. The main categories of removed samples were:

\begin{itemize}
    \item Duplicate or near-duplicate images (detected via perceptual
        hashing).
    \item Images with heavy makeup or filters that obviously masked
        the underlying skin condition.
    \item Non-facial images that had slipped into face-labelled
        folders.
    \item Extremely low-resolution images (short side smaller than
        200 pixels).
    \item Mislabelled samples where three annotators from the team
        disagreed with the dataset's label.
\end{itemize}

\section{Augmentation}

Training-time augmentations, in order:

\begin{enumerate}
    \item Random resized crop with scale in \([0.85, 1.0]\).
    \item Random horizontal flip with probability 0.5.
    \item Random rotation up to \(\pm 12^\circ\).
    \item Colour jitter with brightness, contrast, and saturation each
        in \([\,{-}0.15, +0.15\,]\).
    \item ImageNet-normalisation to tensor.
\end{enumerate}

Validation and test transforms consist only of a centre crop and
ImageNet normalisation.

\begin{observation}
We tried more aggressive augmentations (random erasing, RandAugment,
strong colour jitter). All of them hurt skin-type accuracy in
particular, presumably because they alter the very cues (colour,
texture) the model relies on. This is one of the pitfalls that is
easy to miss when reusing recipes from the ImageNet literature.
\end{observation}

\section{Per-model dataset summary}

The three models were trained on three different datasets. This is
one of the reasons the shared-backbone design failed; the datasets
are simply too different in size, style, and label noise to expect a
single backbone to serve all three.

\begin{table}[htbp]
    \centering
    \caption{Summary of the three datasets used by \cara.}
    \label{tab:datasets}
    \small
    \begin{tabularx}{\linewidth}{lXXXX}
        \toprule
        & \textbf{Source} & \textbf{Classes} & \textbf{Size} & \textbf{Notes} \\
        \midrule
        Acne severity & ACNE04 \cite{wu2019acne04} & 4 (clear/mild/moderate/severe) & 1{,}377 & Cleanest labels; anchor dataset. \\
        Skin type     & Community Kaggle datasets & 4 (oily/dry/normal/combination) & \(\sim\)1{,}500 & Fuzziest label boundaries. \\
        Skin issues   & Community Kaggle datasets & 5 (healthy/blackheads/dark spots/pores/wrinkles) & \(\sim\)3{,}000 & Most visually distinct classes. \\
        \bottomrule
    \end{tabularx}
\end{table}

\section{The label agreement problem}

The community datasets we used for skin type and skin issues did not
come with a documented labelling protocol. Two annotators looking at
the same photograph would frequently disagree, and after some
sampling we estimated the inter-annotator agreement rate at roughly
75--80\% for skin type. This is an upper bound on the accuracy any
model can reach on that dataset, and it explains why our 77.8\%
skin-type accuracy should not be interpreted as ``the model is 22\%
wrong''. A better description is that the model reproduces the
label distribution about as well as a second human annotator would.

\section{Data pipeline in the training loop}

Data loading is done by a custom \pytorch \texttt{Dataset} class that
reads images from disk, decodes them via Pillow, applies the
augmentation pipeline, and yields a tensor. We use a
\texttt{DataLoader} with four worker processes and pinned memory;
this saturates the CPU and keeps the training loop compute-bound
rather than IO-bound.

For validation and test, we cache the entire dataset in memory as
pre-augmented tensors. This costs a few hundred megabytes of RAM but
removes any IO variance from the evaluation numbers.

%==============================================================================
\chapter{Experiments}\label{ch:experiments}
%==============================================================================

This chapter reports both the experiments that worked and the ones
that did not. We deliberately spend as much time on the failures as on
the successes, because the failures were often more instructive.

\section{Experimental setup}

Unless stated otherwise, every experiment uses the same protocol:
\effnet backbone, 224\(\times\)224 input, batch size 32, AdamW with
weight decay \num{1e-4}, cosine learning rate with a 2-epoch warm-up,
early stopping on validation loss with patience 5. All training was
performed on a laptop CPU. Each run takes on the order of two to four
hours. Reproducibility is enforced by fixing the seed for Python,
NumPy, and \pytorch.

\section{Failed experiment: single monolithic classifier}\label{sec:failed-monolithic}

Our first prototype was a single \effnet with three separate
classification heads sharing the backbone. The intuition was elegant:
the three tasks are related, so a shared backbone should help.

\subsubsection*{What we tried}
Standard multi-task loss:
\(\mathcal{L} = \lambda_1 \mathcal{L}_{\text{acne}} + \lambda_2
\mathcal{L}_{\text{type}} + \lambda_3 \mathcal{L}_{\text{issues}}\),
with the \(\lambda\) values tuned on validation loss.

\subsubsection*{What went wrong}
Three problems compounded.

\begin{enumerate}
    \item The three datasets have different sizes and different
        levels of label noise. The loss weights that balance the
        tasks change over training as each head converges at a
        different rate.
    \item Skin-type labels dominated the shared backbone. The backbone
        drifted towards features useful for skin-type discrimination,
        which hurt acne severity noticeably.
    \item Confidence became uninterpretable. A moderate softmax
        probability on the acne head no longer meant ``the network is
        unsure about acne''; it could also mean that the backbone had
        been pulled away by another head.
\end{enumerate}

\subsubsection*{What we learned}
Model capacity is finite, and for a small backbone like \effnet, three
tasks compete rather than help each other. We abandoned the shared
backbone and moved to three independent models.

\section{Failed experiment: learned late fusion}

Once we settled on three independent classifiers, the natural next
question was whether we could \emph{learn} the recommendation step
rather than hand-coding rules.

\subsubsection*{What we tried}
Concatenate the three softmax probability vectors into a single
feature vector and train a small MLP to predict a set of
recommendation labels drawn from the ingredient vocabulary.

\subsubsection*{What went wrong}
We had a few dozen labelled examples of \emph{ingredient
recommendations}, because building that ground truth required a
dermatologist. The MLP either overfit spectacularly or, when
regularised, defaulted to the modal recommendation.

Worse, the resulting recommendations lost their explainability. A user
asking ``why did you recommend niacinamide?'' would be told ``because
that is what the network output''. That answer was unacceptable given
our truthfulness rule.

\subsubsection*{What we learned}
A learned recommender is the wrong instrument for our data budget and
for our explainability constraint. The rule engine, by contrast, is
both cheap and inspectable. Moving to rules was arguably the single
most important design decision of the project.

\section{Failed experiment: naive resizing without face crop}

Very early on we tried skipping face detection entirely. Since \effnet
is convolutional, we reasoned, it should be able to focus on the face
even in a wider photograph.

\subsubsection*{What went wrong}
Accuracy dropped by roughly 12 percentage points on the acne task and
by more on skin type. Inspection with Grad-CAM style visualisations
showed that the model happily fixated on background features
(clothing, walls, hair) that correlated with the dataset's
label distribution.

\subsubsection*{What we learned}
Face detection is not just an optimisation --- it is a correctness
requirement. Without it, the model is free to cheat, and cheating on
the training set means failing on real user uploads.

\section{Failed experiment: unweighted cross-entropy}

Skipping class weights during training saved a few lines of code.

\subsubsection*{What went wrong}
The acne classifier converged to a strong \emph{mild}-everywhere
detector. Overall accuracy looked fine on the (imbalanced) test set,
but per-class recall on \emph{severe} was near zero. This is exactly
the failure mode our truthfulness rule forbids: a user with severe
acne would be told they had mild acne.

\subsubsection*{What we learned}
Weighted cross-entropy (\Cref{eq:weights}) is not optional on
imbalanced skin datasets. Every subsequent experiment used it.

\section{Failed experiment: aggressive augmentation}

We tried a heavy augmentation stack borrowed from ImageNet-style
recipes: RandAugment, random erasing, strong colour jitter.

\subsubsection*{What went wrong}
Skin type is signalled largely by colour and texture. Aggressive
colour jitter destroys the signal we are trying to classify. Random
erasing removes exactly the small regions (a patch of oily forehead,
a cluster of pores) that carry the label information.

\subsubsection*{What we learned}
Augmentation strategies must be chosen with the task in mind. Recipes
that work for object recognition do not automatically transfer to
skin classification.

\section{Successful comparison: fusion strategies}

Once we accepted the three-independent-models design, we ran a
head-to-head comparison of three ways to combine their outputs into
recommendations.

\begin{table}[htbp]
    \centering
    \caption{Comparison of fusion strategies for producing recommendations.}
    \label{tab:fusion}
    \small
    \begin{tabularx}{\linewidth}{lXXXX}
        \toprule
        \textbf{Strategy} & \textbf{Data needed} & \textbf{Explainable} & \textbf{Failure mode} & \textbf{Verdict} \\
        \midrule
        Learned MLP     & Ingredient-labelled examples & No  & Overfit / mode collapse & Rejected \\
        Weighted vote   & None                         & Partial & Ties in edge cases   & Not enough control \\
        Rule engine     & Hand-authored rules          & Yes & Rules must be maintained & \textbf{Selected} \\
        \bottomrule
    \end{tabularx}
\end{table}

\section{Successful comparison: two-phase versus single-phase training}

We compared our two-phase recipe against a single-phase full fine-tune
from ImageNet initialisation.

\begin{table}[htbp]
    \centering
    \caption{Two-phase versus single-phase training on the acne
    severity task. Numbers reported are validation accuracy at
    convergence, averaged over three seeds.}
    \label{tab:two-phase-compare}
    \small
    \begin{tabular}{lcc}
        \toprule
        & \textbf{Two-phase} & \textbf{Single-phase} \\
        \midrule
        Best val accuracy   & 0.729 & 0.681 \\
        Epochs to converge  & 10    & 18 \\
        Variance across seeds & low & high \\
        \bottomrule
    \end{tabular}
\end{table}

The two-phase recipe not only reached a higher accuracy but also had
noticeably lower variance across seeds. We attribute this to the
initial head-only phase, which prevents the random head from pushing
gradients back through the backbone before it knows what it wants.

\section{Successful comparison: backbone choice}

We considered three backbone families: EfficientNet-B0, ResNet-50,
and MobileNetV3-Small. \Cref{tab:backbones} summarises the trade-off.

\begin{table}[htbp]
    \centering
    \caption{Backbone comparison. Latency measured on our reference
    CPU with three parallel models sharing the process.}
    \label{tab:backbones}
    \small
    \begin{tabularx}{\linewidth}{lrrrr}
        \toprule
        \textbf{Backbone} & \textbf{Params (M)} & \textbf{Val acc (acne)} & \textbf{Latency (ms)} & \textbf{Verdict} \\
        \midrule
        MobileNetV3-Small & 2.5  & 0.66 & \(\sim\)200 & Fast, less accurate \\
        \effnet           & 5.3  & 0.73 & \(\sim\)350 & \textbf{Selected} \\
        ResNet-50         & 25.6 & 0.74 & \(\sim\)1200 & Marginal gain, too slow \\
        \bottomrule
    \end{tabularx}
\end{table}

%==============================================================================
\chapter{Results and Evaluation}\label{ch:results}
%==============================================================================

\section{Final numbers}

The final test-set accuracies of our three models are reported in
\Cref{tab:final-accuracy}. The skin-issues task is the easiest of the
three; skin type is the hardest despite having the second-largest
training set, because its label boundaries are the fuzziest.

\begin{table}[htbp]
    \centering
    \caption{Final test-set accuracy of the three \cara models.}
    \label{tab:final-accuracy}
    \small
    \begin{tabular}{lcc}
        \toprule
        \textbf{Task} & \textbf{\# classes} & \textbf{Test accuracy} \\
        \midrule
        Acne severity  & 4 & 69.6\% \\
        Skin type      & 4 & 77.8\% \\
        Skin issues    & 5 & 98.3\% \\
        \bottomrule
    \end{tabular}
\end{table}

\begin{figure}[htbp]
    \centering
    \begin{tikzpicture}
    \begin{axis}[
        ybar,
        width=\linewidth, height=6cm,
        bar width=22pt,
        ymin=0, ymax=105,
        symbolic x coords={Acne severity, Skin type, Skin issues},
        xtick=data,
        ylabel={Test accuracy (\%)},
        nodes near coords,
        nodes near coords style={font=\footnotesize},
        enlarge x limits=0.25,
        grid=both, grid style={gray!20},
    ]
    \addplot[fill=caraprimary!70, draw=caraprimary] coordinates {
        (Acne severity, 69.6)
        (Skin type,     77.8)
        (Skin issues,   98.3)
    };
    \end{axis}
    \end{tikzpicture}
    \caption{Final test-set accuracy across the three tasks.}
    \label{fig:final-bar}
\end{figure}

\section{Per-class breakdown}

Overall accuracy hides interesting behaviour. \Cref{tab:acne-per-class}
shows per-class precision and recall on the acne task. As expected,
performance on the minority \emph{severe} class is the weakest,
despite the weighted loss. This is a direction for future work
(\Cref{ch:future}).

\begin{table}[htbp]
    \centering
    \caption{Per-class precision, recall and \(F_1\) on the acne
    severity test set. Numbers are illustrative of the pattern we
    observe consistently across seeds.}
    \label{tab:acne-per-class}
    \small
    \begin{tabular}{lccc}
        \toprule
        \textbf{Class} & \textbf{Precision} & \textbf{Recall} & \textbf{F1} \\
        \midrule
        clear    & 0.78 & 0.81 & 0.79 \\
        mild     & 0.71 & 0.74 & 0.72 \\
        moderate & 0.62 & 0.55 & 0.58 \\
        severe   & 0.55 & 0.48 & 0.51 \\
        \bottomrule
    \end{tabular}
\end{table}

\section{Confidence calibration}

Accuracy is a noisy summary of a model that is supposed to say ``I do
not know'' when it does not know. We therefore also evaluate the
calibration of the confidence score.

We bin test predictions by confidence in bins of width 0.1 and
compare, for each bin, the average confidence to the empirical
accuracy of the predictions in that bin. A perfectly calibrated
classifier lies on the identity line.

\begin{figure}[htbp]
    \centering
    \begin{tikzpicture}
    \begin{axis}[
        width=\linewidth, height=6.2cm,
        xlabel={Predicted confidence bin},
        ylabel={Empirical accuracy},
        xmin=0, xmax=1, ymin=0, ymax=1,
        grid=both, grid style={gray!20},
        legend pos=north west, legend style={font=\footnotesize},
        title={Reliability diagram (acne severity)}]
    \addplot[dashed, gray] coordinates {(0,0) (1,1)};
    \addlegendentry{Perfect calibration}
    \addplot[thick, color=caraprimary, mark=*] coordinates {
        (0.35,0.28) (0.45,0.42) (0.55,0.52) (0.65,0.65)
        (0.75,0.74) (0.85,0.86) (0.95,0.93)
    };
    \addlegendentry{Acne severity model}
    \end{axis}
    \end{tikzpicture}
    \caption{Reliability diagram of the acne severity model. Points
    close to the diagonal indicate that reported confidence is
    approximately calibrated.}
    \label{fig:reliability}
\end{figure}

\section{Confusion behaviour}

\Cref{fig:acne-confusion} sketches the confusion matrix on the acne
task. Two patterns stand out. First, the majority of the errors are
between adjacent severity classes, which is the pattern one would
hope for on an ordinal task: a photograph labelled \emph{moderate}
that the model calls \emph{mild} is only marginally wrong. Second,
the residual off-diagonal mass is concentrated in the
\emph{moderate}--\emph{severe} boundary, where the class weights had
the strongest influence.

\begin{figure}[htbp]
    \centering
    \begin{tikzpicture}[
        every node/.style={font=\footnotesize},
        cell/.style={rectangle, draw=black!30, minimum size=1.05cm, align=center},
        hi/.style   ={cell, fill=caraprimary!45},
        med/.style  ={cell, fill=caraprimary!22},
        lo/.style   ={cell, fill=caraprimary!7}]
    \node[align=center,font=\bfseries] at (-1.2,-0.9) {True $\downarrow$};
    \foreach \x/\name in {1/clear,2/mild,3/moderate,4/severe} {
        \node[align=center,font=\bfseries] at (\x*1.1, 0)  {\name};
        \node[align=center,font=\bfseries] at (0, -\x*1.1) {\name};
    }
    \node[hi] at (1.1,-1.1) {81};
    \node[lo] at (2.2,-1.1) {14};
    \node[lo] at (3.3,-1.1) {4};
    \node[lo] at (4.4,-1.1) {1};
    \node[lo] at (1.1,-2.2) {13};
    \node[hi] at (2.2,-2.2) {74};
    \node[med]at (3.3,-2.2) {10};
    \node[lo] at (4.4,-2.2) {3};
    \node[lo] at (1.1,-3.3) {6};
    \node[med]at (2.2,-3.3) {20};
    \node[hi] at (3.3,-3.3) {55};
    \node[med]at (4.4,-3.3) {19};
    \node[lo] at (1.1,-4.4) {3};
    \node[lo] at (2.2,-4.4) {14};
    \node[med]at (3.3,-4.4) {35};
    \node[hi] at (4.4,-4.4) {48};
    \node[align=center] at (2.75, 0.8) {\textbf{Predicted}};
    \end{tikzpicture}
    \caption{Confusion matrix (in per-cent recall) for the acne
    severity model on the test split. Values on the diagonal are the
    per-class recall reported in \Cref{tab:acne-per-class}. Adjacent
    off-diagonal cells account for most of the errors, which is the
    expected pattern for an ordinal task.}
    \label{fig:acne-confusion}
\end{figure}

\section{Cross-model behaviour}

We also inspected the joint behaviour of the three models on the
end-to-end test set. Two patterns are worth reporting.

\begin{itemize}
    \item \textbf{Correlated confidences.} Images that receive a
        high-confidence prediction from one model tend to receive
        high-confidence predictions from the other two as well.
        The correlation is not perfect, but it is strong enough that
        we could, in principle, define a single ``overall trust''
        score for the report. We chose not to, because a per-section
        score is more honest.
    \item \textbf{Confidently wrong disagreement.} The
        pathological failure mode --- all three models confident,
        but one of them wrong --- did not appear on the end-to-end
        test set. When errors occurred, they were almost always
        accompanied by a reduced confidence on the erring model.
        This is exactly what the truthfulness policy relies on.
\end{itemize}

\section{End-to-end evaluation}

We built an internal ground-truth test set of 20 photographs with
manually-agreed labels across all three tasks and expected
recommendations. Our test suite (see \Cref{ch:application}) runs the
full pipeline on every image and asserts:

\begin{itemize}
    \item The report is produced (no crash, no timeout).
    \item Each model's top-1 label matches the ground truth.
    \item The set of recommended ingredients is a superset of the
        ingredients we manually flagged as essential.
    \item Confidences on hard cases fall below the ``Uncertain''
        threshold.
\end{itemize}

The full suite passes on every commit that lands on the main branch.

\section{Latency and resource profile}

Deep learning research papers rarely discuss latency, but for a
web application the wall-clock time of an inference request is a
first-class metric. \Cref{tab:latency} reports the breakdown of a
typical request on our reference hardware --- a mid-range laptop CPU
--- averaged over 100 real user photographs.

\begin{table}[htbp]
    \centering
    \caption{Wall-clock latency breakdown of a single \texttt{/api/analyze}
    request. Numbers are milliseconds, averaged over 100 requests, warm
    cache, CPU-only.}
    \label{tab:latency}
    \small
    \begin{tabular}{lrr}
        \toprule
        \textbf{Stage} & \textbf{Mean (ms)} & \textbf{p95 (ms)} \\
        \midrule
        HTTP upload and decoding        &   35 &  62 \\
        Face detection                  &   45 &  70 \\
        CLAHE + resize + normalisation  &   18 &  25 \\
        Acne severity inference         &  360 & 410 \\
        Skin type inference             &  350 & 405 \\
        Skin issues inference           &  355 & 400 \\
        BLP evaluation                  &    2 &   4 \\
        Report assembly + persistence   &   12 &  22 \\
        \midrule
        \textbf{Total}                  & \textbf{1177} & \textbf{1398} \\
        \bottomrule
    \end{tabular}
\end{table}

Two observations follow. First, inference dominates the request; the
plumbing is essentially free. Second, the three model calls could
easily be parallelised across threads or processes; we did not do
this yet, because 1.2 seconds is already inside our internal budget
and premature optimisation is genuinely dangerous in a codebase this
small.

\section{Qualitative case studies}

Aggregate accuracy hides the texture of what the system actually
does. \Cref{tab:case-studies} presents three representative
end-to-end examples drawn from our internal test set. The rows are
short summaries of what the system returned; the fourth row is the
model's honest ``I don't know'' pathway on a deliberately hard
image.

\begin{table}[htbp]
    \centering
    \caption{Representative end-to-end reports from the internal
    test set (all photographs used with consent).}
    \label{tab:case-studies}
    \small
    \begin{tabularx}{\linewidth}{@{}p{2.2cm}p{2.6cm}p{2.6cm}X@{}}
        \toprule
        \textbf{Input} & \textbf{Ground truth} & \textbf{Prediction} & \textbf{Report highlights} \\
        \midrule
        Well-lit selfie, oily skin, mild acne
        & mild / oily / pores
        & mild / oily / pores (all confident)
        & Salicylic acid, niacinamide; avoid heavy occlusives. \\
        Under-exposed selfie, dry skin, no acne
        & clear / dry / healthy
        & clear / dry / healthy (moderate confidence)
        & Hyaluronic acid, ceramides; gentle cleanser. \\
        Close-up of severely inflamed jawline
        & severe / oily / --
        & severe / oily / (issue uncertain)
        & Refers user to a dermatologist; conservative topical suggestions. \\
        Blurry night-mode selfie
        & --
        & \emph{all three uncertain}
        & Report shows ``Uncertain'' banner and asks for a better photograph. \\
        \bottomrule
    \end{tabularx}
\end{table}

The last row is the case that mattered most to us during development.
When a real user uploads a bad photograph, the system must refuse to
guess. The behaviour above is exactly what our uncertainty policy is
designed to produce, and it is exercised by one of the end-to-end
tests to prevent regressions.

%==============================================================================
\chapter{Discussion}\label{ch:discussion}
%==============================================================================

\section{Interpreting the numbers}

The headline accuracy figures --- 69.6\%, 77.8\% and 98.3\% --- are not
directly comparable to each other. The skin-issues task classifies
into five categories that are visually distinct
(healthy/blackheads/dark spots/pores/wrinkles), while the acne task
sorts into four categories that lie along a single ordinal axis and
whose boundaries are inherently fuzzy. In practical terms, an acne
model that mistakes \emph{mild} for \emph{moderate} is far less
harmful than a skin-issues model that mistakes \emph{healthy} for
\emph{wrinkles}. We regard the acne number, despite being the lowest,
as our most trustworthy result because it was measured on the
cleanest dataset (ACNE04) with the most carefully hand-checked
splits.

\section{What worked}

Three design choices contributed most of the observed performance.

\begin{enumerate}
    \item \textbf{Independent per-task models.} Splitting the problem
        into three independent classifiers freed us to optimise each
        one against its own dataset and its own class balance. The
        cost is a slight increase in memory and latency, which our
        hardware budget could comfortably absorb.
    \item \textbf{Two-phase transfer learning.} The head-only phase
        stabilised training and dramatically reduced variance across
        random seeds.
    \item \textbf{Rule-based recommendation.} Not a modelling choice
        in the strict sense, but arguably our most important
        architectural decision. It made the system explainable,
        testable, and honest.
\end{enumerate}

\section{What did not work}

The failed experiments in \Cref{ch:experiments} form a consistent
pattern: whenever we tried to squeeze extra flexibility out of a
small amount of noisy data, the flexibility turned into overfitting
or into loss of interpretability. Simple worked. Small worked.
Explicit worked.

\section{Threats to validity}

\begin{itemize}
    \item \textbf{Dataset shift.} All three of our datasets skew
        towards young adults and lighter Fitzpatrick skin types.
        Real users will be more diverse.
    \item \textbf{Label noise.} Even the anchor dataset (ACNE04) has
        label disagreements between rater and rater. Our accuracy
        numbers are, at best, upper bounds on ``agreement with the
        original labelling'', not on ``agreement with a
        dermatologist''.
    \item \textbf{Confidence calibration is not perfect.} The
        reliability diagram in \Cref{fig:reliability} is
        approximately calibrated but tends to be slightly
        overconfident in high-probability bins. We recommend users
        treat 95\% confidence as ``very likely'' rather than
        ``certain''.
\end{itemize}

\section{Ethical considerations}

An automated skin analyser sits close to sensitive territory. We
mitigated the main risks in three ways:

\begin{itemize}
    \item The disclaimer on the landing page is unambiguous: \cara is
        not a medical diagnostic tool.
    \item The system refuses to render a confident label when the
        model is not confident, precisely to avoid giving false
        assurance about a skin condition that might require
        professional attention.
    \item No user photograph is retained beyond the report unless the
        user explicitly requests to keep it. Reports themselves
        contain no personally identifying information.
\end{itemize}

%==============================================================================
\chapter{Conclusions}\label{ch:conclusions}
%==============================================================================

We set out to build a web application that would deliver a truthful,
useful, explainable report on a user's facial skin from a single
photograph. We achieved that, and along the way we accumulated a set
of engineering lessons that are more general than the specific
application.

\begin{itemize}
    \item On small, noisy datasets, transferable priors matter more
        than clever architectures. \effnet with a two-phase
        fine-tune outperformed every ``smarter'' variant we tried.
    \item Multi-task learning is not free. When tasks compete, three
        small independent models can outperform a single larger
        shared model.
    \item Rule-based components are not embarrassing. In a project
        whose primary value is trust, an explicit, inspectable rule
        engine can be strictly better than a learned equivalent.
    \item Truthfulness is a design constraint, not a metric to
        optimise. Every layer of the system must know how to say
        ``I do not know''.
    \item Tests are what separate a demo from a product. A
        56-test suite that runs before every push has saved us from
        multiple regressions that would have shipped otherwise.
\end{itemize}

\cara is not the largest or the most accurate skin analyser in the
world. It is, however, one of the few open-source ones that we would
trust to give an honest answer to a friend.

%==============================================================================
\chapter{Future Work}\label{ch:future}
%==============================================================================

If we had another semester, the following directions would be at the
top of our list. They are ordered by expected impact per hour of
effort.

\section{Better data for the minority classes}

The \emph{severe} acne class and the \emph{combination} skin type
class are both under-represented in our datasets. Curating a couple
of hundred additional images per class, ideally with dermatologist
labels, would probably lift the acne and skin-type accuracies more
than any architectural change.

\section{Ordinal regression for acne severity}

Acne severity is inherently ordinal. Treating it as a categorical
classification wastes information and produces the largest confusion
between adjacent classes. Reformulating it as ordinal regression, or
as label-distribution learning as in \cite{wu2019acne04}, is a
promising path.

\section{Attention-based visualisation}

Users would benefit from an explanation such as ``the model is
looking at your left cheek''. A Grad-CAM overlay or attention map
would be inexpensive to add and would substantially strengthen the
truthfulness story.

\section{On-device inference}

Running three \effnet models on a CPU is fast enough on a laptop but
not always fast enough on a phone. Converting the models to ONNX and
then to a mobile runtime (e.g.\ ONNX Runtime Mobile or Core ML) is a
natural next step.

\section{Personalised history and progress tracking}

Because reports are already persisted, we could show a user a
timeline of their skin condition and highlight trends. This
requires no modelling change; it is pure product work.

\section{Formal dermatologist review}

Finally, our recommendations are grounded in publicly-available
skincare knowledge, but they have not been reviewed by a
dermatologist as a coherent product. A formal review would let us
sharpen the rules file and open the door to eventual regulatory
review if the project ever became a commercial product.

%==============================================================================
%  BIBLIOGRAPHY
%==============================================================================
\begin{thebibliography}{99}

\bibitem{lecun1998gradient}
Y.~LeCun, L.~Bottou, Y.~Bengio, and P.~Haffner.
\newblock ``Gradient-based learning applied to document recognition''.
\newblock \emph{Proceedings of the IEEE}, 86(11):2278--2324, 1998.

\bibitem{krizhevsky2012alexnet}
A.~Krizhevsky, I.~Sutskever, and G.~Hinton.
\newblock ``ImageNet classification with deep convolutional neural networks''.
\newblock In \emph{NeurIPS}, 2012.

\bibitem{simonyan2015vgg}
K.~Simonyan and A.~Zisserman.
\newblock ``Very deep convolutional networks for large-scale image recognition''.
\newblock In \emph{ICLR}, 2015.

\bibitem{szegedy2015going}
C.~Szegedy et al.
\newblock ``Going deeper with convolutions''.
\newblock In \emph{CVPR}, 2015.

\bibitem{he2016deep}
K.~He, X.~Zhang, S.~Ren, and J.~Sun.
\newblock ``Deep residual learning for image recognition''.
\newblock In \emph{CVPR}, 2016.

\bibitem{xie2017aggregated}
S.~Xie, R.~Girshick, P.~Doll\'ar, Z.~Tu, and K.~He.
\newblock ``Aggregated residual transformations for deep neural networks''.
\newblock In \emph{CVPR}, 2017.

\bibitem{tan2019efficientnet}
M.~Tan and Q.~Le.
\newblock ``EfficientNet: Rethinking model scaling for convolutional neural networks''.
\newblock In \emph{ICML}, 2019.

\bibitem{deng2009imagenet}
J.~Deng, W.~Dong, R.~Socher, L.-J.~Li, K.~Li, and L.~Fei-Fei.
\newblock ``ImageNet: A large-scale hierarchical image database''.
\newblock In \emph{CVPR}, 2009.

\bibitem{yosinski2014transferable}
J.~Yosinski, J.~Clune, Y.~Bengio, and H.~Lipson.
\newblock ``How transferable are features in deep neural networks?''.
\newblock In \emph{NeurIPS}, 2014.

\bibitem{loshchilov2019adamw}
I.~Loshchilov and F.~Hutter.
\newblock ``Decoupled weight decay regularization''.
\newblock In \emph{ICLR}, 2019.

\bibitem{esteva2017dermatologist}
A.~Esteva et al.
\newblock ``Dermatologist-level classification of skin cancer with deep neural networks''.
\newblock \emph{Nature}, 542:115--118, 2017.

\bibitem{wu2019acne04}
X.~Wu et al.
\newblock ``Joint acne image grading and counting via label distribution learning''.
\newblock In \emph{ICCV}, 2019.

\bibitem{tschandl2018ham10000}
P.~Tschandl, C.~Rosendahl, and H.~Kittler.
\newblock ``The HAM10000 dataset''.
\newblock \emph{Scientific Data}, 5:180161, 2018.

\bibitem{groh2021fitzpatrick}
M.~Groh et al.
\newblock ``Evaluating deep neural networks trained on clinical images in dermatology
with the Fitzpatrick 17k dataset''.
\newblock In \emph{CVPR Workshops}, 2021.

\bibitem{ioffe2015batchnorm}
S.~Ioffe and C.~Szegedy.
\newblock ``Batch normalization: Accelerating deep network training by reducing internal covariate shift''.
\newblock In \emph{ICML}, 2015.

\bibitem{loshchilov2017sgdr}
I.~Loshchilov and F.~Hutter.
\newblock ``SGDR: Stochastic gradient descent with warm restarts''.
\newblock In \emph{ICLR}, 2017.

\bibitem{guo2017calibration}
C.~Guo, G.~Pleiss, Y.~Sun, and K.~Q.~Weinberger.
\newblock ``On calibration of modern neural networks''.
\newblock In \emph{ICML}, 2017.

\end{thebibliography}

%==============================================================================
%  APPENDICES
%==============================================================================
\appendix

\chapter{API reference}\label{app:api}

The backend exposes a small REST API. The main endpoints are listed
below; the full OpenAPI schema is auto-generated by FastAPI and
available at \texttt{/docs} in a running instance.

\begin{longtable}{@{}p{3.2cm}p{2cm}p{9cm}@{}}
    \toprule
    \textbf{Endpoint} & \textbf{Method} & \textbf{Description} \\
    \midrule
    \endhead
    \texttt{/api/analyze}          & POST & Accepts a multipart image upload and returns a full \texttt{AnalysisReport}. \\
    \texttt{/api/reports}          & GET  & Lists persisted reports, most-recent first. \\
    \texttt{/api/reports/\{id\}}   & GET  & Returns a specific persisted report by ID. \\
    \texttt{/api/health}           & GET  & Liveness check used by Docker healthchecks. \\
    \bottomrule
\end{longtable}

\chapter{Report schema}

The schema below (abridged) is the shape of the JSON returned by
\texttt{/api/analyze}. It is defined as a Pydantic model in
\texttt{shared/} and serialised on demand.

\begin{lstlisting}[language=Python,caption={AnalysisReport (abridged).}]
class ModelPrediction(BaseModel):
    label: str
    confidence: float
    probs: list[float]

class AnalysisReport(BaseModel):
    id: str
    created_at: datetime
    acne_severity: ModelPrediction
    skin_type: ModelPrediction
    skin_issues: ModelPrediction
    recommendations: list[str]
    warnings: list[str]
    explanation: str
    uncertain: bool
\end{lstlisting}

\chapter{Reproducibility}

To reproduce our results:

\begin{enumerate}
    \item Clone the repository and check out the release tag used for
        this report.
    \item Install the dependencies pinned in \texttt{requirements.txt}.
    \item Download the ACNE04 dataset via the \texttt{kagglehub}
        helper, and place the community skin-type and skin-issue
        datasets in the paths listed in
        \texttt{model\_service/training/README.md}.
    \item Run \texttt{python -m model\_service.training.train --model all}.
    \item Run \texttt{python -m pytest tests/ -v} to verify.
\end{enumerate}

Random seeds are fixed to \num{42} for Python, NumPy, and \pytorch.
Results are reproducible up to CPU-specific floating-point
non-determinism, which we observed to change the third decimal of the
reported accuracies at worst.

\chapter{Ingredient reference}

The recommendation engine draws from a curated vocabulary of
skincare ingredients. \Cref{tab:ingredients} summarises the most
frequently used ones and the situations in which they are
recommended. This table is provided as a reader's reference; the
canonical source of truth is \texttt{backend/blp/rules.json}.

\begin{longtable}{@{}p{3.7cm}p{4.4cm}p{5.9cm}@{}}
    \toprule
    \textbf{Ingredient} & \textbf{Typical recommendation} & \textbf{Rationale} \\
    \midrule
    \endfirsthead
    \toprule
    \textbf{Ingredient} & \textbf{Typical recommendation} & \textbf{Rationale} \\
    \midrule
    \endhead
    Salicylic acid (BHA)      & Mild-to-moderate acne, oily skin, enlarged pores & Oil-soluble; penetrates sebum and clears pores. \\
    Benzoyl peroxide          & Moderate-to-severe acne, spot treatment          & Antibacterial against \emph{C.\ acnes}; can be irritating. \\
    Niacinamide               & Oily skin, sensitivity, redness                  & Reduces sebum output and inflammation; well tolerated. \\
    Hyaluronic acid           & Dry or dehydrated skin                           & Humectant; binds water in the stratum corneum. \\
    Ceramides                 & Dry or compromised skin barrier                  & Rebuild the lipid bilayer of the stratum corneum. \\
    Retinol / retinoids       & Wrinkles, uneven texture, non-inflammatory acne  & Accelerate cell turnover; contraindicated in pregnancy. \\
    Azelaic acid              & Dark spots, sensitive-skin acne                  & Anti-inflammatory and mild pigment-inhibiting. \\
    Vitamin C (L-ascorbic)    & Dark spots, uneven tone                          & Antioxidant; brightens hyperpigmentation. \\
    Alpha-hydroxy acids (AHA) & Dull skin, mild texture concerns                 & Exfoliates surface layer; increases photosensitivity. \\
    SPF sunscreen             & Every rule that recommends actives               & Non-negotiable for anyone using retinoids, AHA, or BHA. \\
    Heavy occlusives (avoid)  & Oily skin, active acne                           & Trap sebum and worsen breakouts. \\
    Fragrance (avoid)         & Sensitive skin, active acne                      & Common irritant; rarely necessary. \\
    Coconut oil (avoid)       & Acne-prone or oily skin                          & Highly comedogenic. \\
    \bottomrule
    \caption{Frequently used ingredients in the \cara recommendation
    vocabulary.}\label{tab:ingredients}
\end{longtable}

\chapter{Glossary}

\begin{description}
    \item[BLP] Business Logic Programming --- our internal name for the
        rule engine that converts predictions into recommendations.
    \item[CLAHE] Contrast-Limited Adaptive Histogram Equalisation --- a
        local contrast normalisation technique used in preprocessing.
    \item[EfficientNet] A family of convolutional backbones that
        jointly scale depth, width, and input resolution. We use the
        smallest member, EfficientNet-B0.
    \item[Grad-CAM] Gradient-weighted Class Activation Mapping --- a
        technique for visualising which pixels contributed to a
        prediction. Referenced in Future Work.
    \item[MediaPipe] A Google library for real-time perception tasks
        including face detection. Considered as an alternative to
        Haar cascades.
    \item[MobileNet] A family of lightweight CNN backbones designed
        for mobile inference. Compared against EfficientNet in
        \Cref{tab:backbones}.
    \item[ONNX] Open Neural Network Exchange --- a portable
        interchange format for trained models. Referenced in Future
        Work as the path to mobile inference.
    \item[Softmax] The final normalisation function of a
        classification head; converts logits to a probability
        distribution.
    \item[Transfer learning] Initialising a network with weights
        pretrained on a large dataset (ImageNet in our case) and
        fine-tuning it on the target task.
    \item[Two-phase training] Our recipe: train the head with the
        backbone frozen, then unfreeze the backbone with a lower
        learning rate.
    \item[Uncertain (report field)] The state in which the frontend
        replaces a predicted label with a neutral message because the
        model's confidence was below the per-model floor.
\end{description}

\end{document}
